\documentclass[a4paper,12pt]{article}
\usepackage[T2A]{fontenc}
\usepackage[utf8]{inputenc}
\usepackage[russian]{babel}
\usepackage{amsmath,amssymb,mathtools}
\usepackage{DejaVuSans}
\renewcommand{\familydefault}{\sfdefault}
\usepackage{geometry}
\geometry{left=24mm,right=24mm,top=23mm,bottom=24mm}
\usepackage{microtype}
\usepackage{icomma}
\usepackage{booktabs,tabularx,array}
\usepackage{enumitem}
\usepackage{xcolor}
\usepackage{tikz}
\usetikzlibrary{arrows.meta,positioning,shapes.geometric,calc}
\usepackage{fancyhdr}
\usepackage{setspace}
\usepackage{hyperref}
\hypersetup{colorlinks=true,linkcolor=teal!65!black,urlcolor=teal!65!black}
\definecolor{accent}{HTML}{087E82}
\definecolor{deep}{HTML}{075B60}
\definecolor{ink}{HTML}{292D31}
\definecolor{muted}{HTML}{617075}
\color{ink}
\setstretch{1.22}
\setlength{\parindent}{0pt}
\setlength{\parskip}{5pt}
\setlist[itemize]{leftmargin=7mm,itemsep=2pt,topsep=2pt}
\setlist[enumerate]{leftmargin=8mm,itemsep=5pt,topsep=5pt}
\setlength{\headheight}{15pt}
\pagestyle{fancy}
\fancyhf{}
\fancyhead[L]{\small\color{muted}Дробь от числа, проценты и задачи}
\fancyhead[R]{\small\color{muted}Надежда Клименко}
\fancyfoot[C]{\small\color{muted}\thepage}
\renewcommand{\headrulewidth}{0.35pt}
\renewcommand{\footrulewidth}{0pt}
\newcommand{\ruleline}{\par\vspace{-2pt}{\color{accent}\hrule height 0.55pt}\vspace{8pt}}
\newcommand{\concept}[1]{\par\smallskip{\color{deep}\bfseries #1}\par\smallskip}
\newcommand{\method}[1]{{\color{accent}\bfseries #1}}
\newcommand{\example}[1]{\par\medskip{\color{deep}\textbf{Пример.} #1}\par\smallskip}
\newcommand{\checkline}{\par\smallskip{\color{accent}\rule{14mm}{0.7pt}}\par\smallskip}
\newcommand{\mysection}[2]{\section*{\color{deep}#1. #2}\addcontentsline{toc}{section}{#1. #2}\ruleline}
\begin{document}
\begin{titlepage}
\thispagestyle{empty}
\begin{center}
\vspace*{30mm}
{\large\color{deep}\textbf{МАТЕМАТИКА}}\par
\vspace{19mm}
{\LARGE Чек-лист}\par
\vspace{12mm}
{\LARGE\color{accent}Дробь от числа, проценты\\[2mm]и задачи на составление уравнений}\par
\vspace{20mm}
{\large Повторение приёмов занятия}\par
\vfill
{\large Надежда Клименко}\par
\vspace{12mm}
{\normalsize 3 ноября 2026 г.}\par
\vspace{23mm}
{\small Лёвин Артём Александрович\\[1mm]эксклюзивно для Надежды Клименко}\par
\vspace{17mm}
\begin{tikzpicture}[x=1cm,y=1cm]
\draw[accent,line width=0.8pt] (0,0) .. controls (1.4,0.62) and (2.4,-0.52) .. (3.6,0)
 .. controls (4.8,0.60) and (5.8,-0.52) .. (7.2,0);
\end{tikzpicture}
\end{center}
\end{titlepage}

\mysection{1}{Дробь от числа и число по его доле}
\concept{Что обозначает дробь}
Если целое разделили на $b$ равных частей и взяли $a$ из них, получили долю $\dfrac ab$. Например, $\dfrac38$ числа --- три из восьми равных частей этого числа.

\concept{Если известно целое, ищем его часть}
\[\boxed{\text{Часть}=\text{целое}\cdot\text{доля}}\qquad A=C\cdot\frac ab.\]
\example{Найдите $\dfrac35$ от $140$.}
\[140\cdot\frac35=\frac{140}{5}\cdot3=28\cdot3=84.\]

\concept{Если известна часть, ищем целое}
\[\boxed{\text{Целое}=\text{часть}:\text{доля}}\qquad C=A:\frac ab\quad(a\ne0).\]
\example{Число $54$ составляет $\dfrac3{10}$ неизвестного числа. Найдите его.}
\[54:\frac3{10}=54\cdot\frac{10}{3}=180.\]
\method{Самопроверка.} Если доля меньше единицы, то найденная часть меньше исходного целого. Если ищем целое по его положительной правильной доле, ответ больше известной части.

\mysection{2}{Проценты как дроби}
Процент --- одна сотая часть целого: $1\%=\dfrac1{100}=0{,}01$. Всегда выясняйте, \textbf{от какой величины} вычисляются проценты.
\[15\%=0{,}15,\qquad 24\%=0{,}24,\qquad 60\%=0{,}6,\qquad 95\%=0{,}95.\]

\concept{Найти первоначальное целое по известной доле}
\example{Убрали $3600$ гектаров --- это $60\%$ всей площади. Найдите всю площадь.}
\[0{,}6x=3600,\qquad x=\frac{3600}{0{,}6}=6000\ \text{га}.\]
\method{Проверка:} $6000\cdot0{,}6=3600$. Не путайте деление на $0{,}6$ с умножением на $0{,}6$.

\concept{На сколько изменилась величина}
При увеличении на $p\%$ новая величина равна
\[C\left(1+\frac p{100}\right).\]
При уменьшении на $p\%$ новая величина равна
\[C\left(1-\frac p{100}\right).\]
А \textbf{само изменение} равно $C\cdot\dfrac p{100}$. Именно эта формула нужна, когда в условии сказано «на столько-то больше» или «уменьшилось на столько-то».

\example{Станок стал изготавливать на $105$ деталей больше после повышения производительности на $15\%$. Сколько деталей изготавливал раньше?}
\[0{,}15x=105,\qquad x=\frac{105}{0{,}15}=700.\]
\method{Проверка:} $15\%$ от $700$ --- это $105$, значит новая выработка $805$ деталей.

\example{Товар подешевел на $24\%$, то есть на $57{,}6$ руб. Сколько он стоил до скидки?}
\[0{,}24x=57{,}6,\qquad x=\frac{57{,}6}{0{,}24}=240\ \text{руб.}\]
Цена после скидки: $240\cdot0{,}76=182{,}4$ руб.

\clearpage
\mysection{3}{Десятичные и смешанные дроби в уравнениях}
\concept{Деление десятичных дробей}
Можно умножить делимое и делитель на одно и то же число, отличное от нуля, чтобы в делителе исчезла запятая:
\[0{,}16:0{,}5=1{,}6:5=0{,}32.\]
Или перейти к обыкновенным дробям:
\[\frac{0{,}16}{0{,}5}=\frac{16/100}{5/10}=\frac{16}{50}=\frac8{25}=0{,}32.\]
\method{Важно.} Изменяя делитель в $10$ или $100$ раз, изменяйте во столько же раз делимое.

\concept{Смешанное число в роли коэффициента}
\[5\frac14=\frac{21}{4}=5{,}25.\]
\example{Решите уравнение $5\dfrac14y-5\dfrac14=5\dfrac14$.}
При переносе слагаемого из левой части в правую его знак меняется:
\begin{align*}
5\frac14y&=5\frac14+5\frac14,\\
5\frac14y&=2\cdot5\frac14,\\
y&=\frac{2\cdot5\frac14}{5\frac14}=2.
\end{align*}
\method{Проверка:} $5\dfrac14\cdot2-5\dfrac14=5\dfrac14$.

\mysection{4}{Отношение величин и уравнение по условию}
Если одна величина составляет $95\%$ другой, то при обозначении второй через $x$ первая равна $0{,}95x$. Если она в $1\dfrac25$ раза больше, то равна $\dfrac75x$.

\example{Два отряда посадили $780$ саженцев. Первый посадил $95\%$ от количества второго. Сколько посадил каждый?}
Пусть второй отряд посадил $x$ саженцев. Тогда первый --- $0{,}95x$.
\[x+0{,}95x=780,\qquad 1{,}95x=780,\qquad x=400.\]
Следовательно, первый отряд посадил $380$ саженцев, второй --- $400$.
\method{Проверка:} $380+400=780$ и $380:400=0{,}95$.

\example{Смежные углы составляют $180^\circ$. Один угол в $1\dfrac25$ раза больше другого. Найдите оба.}
Обозначим меньший угол через $x$. Больший угол равен $\dfrac75x$.
\[x+\frac75x=180^\circ,\qquad \frac{12}{5}x=180^\circ,\qquad x=75^\circ.\]
Больший угол: $75^\circ\cdot\dfrac75=105^\circ$.
\method{Проверка:} $75^\circ+105^\circ=180^\circ$.

\clearpage
\mysection{5}{Задачи на движение навстречу друг другу}
Если два объекта движутся навстречу, расстояние между ними сокращается со \textbf{скоростью сближения}
\[v_{\text{сбл}}=v_1+v_2.\]
Для времени до встречи $t$ и начального расстояния $S$:
\[\boxed{S=(v_1+v_2)t}.\]
При движении в одном направлении скорость сближения при догонянии равна разности скоростей: $v_{\text{сбл}}=v_{\text{быстр}}-v_{\text{медл}}$.

\example{Автомобили находятся в $126$ км друг от друга и едут навстречу. Они встретятся через $\dfrac{14}{15}$ ч. Скорость одного автомобиля составляет $80\%$ скорости другого. Найдите скорости.}
Пусть скорость более быстрого автомобиля $x$ км/ч. Скорость второго --- $0{,}8x$ км/ч. Тогда
\begin{align*}
(x+0{,}8x)\cdot\frac{14}{15}&=126,\\
1{,}8x&=126:\frac{14}{15}=135,\\
x&=75,\qquad 0{,}8x=60.
\end{align*}
\method{Проверка:} $(75+60)\cdot\dfrac{14}{15}=126$ км.

\mysection{6}{Как проверять составленное уравнение}
\begin{enumerate}
\item \textbf{Назовите целое.} Укажите, относительно какой величины берут дробь или проценты.
\item \textbf{Введите переменную.} Напишите, что именно означает $x$ и в каких единицах.
\item \textbf{Выразите остальные величины через $x$.} Например, $95\%$ от $x$ --- это $0{,}95x$.
\item \textbf{Запишите связь из условия.} Сумма, разность, доля, путь или изменение дают уравнение.
\item \textbf{Решите и проверьте.} Подставьте ответ в исходное условие, оцените его смысл и единицы измерения.
\end{enumerate}
\concept{Частые ошибки}
\begin{itemize}
\item При поиске целого по известной доле умножают на долю вместо деления.
\item При увеличении на $15\%$ путают прибавку $0{,}15x$ с новой величиной $1{,}15x$.
\item Считают, что $95\%$ «от второго» означает $0{,}95$ от общей суммы.
\item В уравнении меняют знак слагаемого, хотя ещё не перенесли его через знак равенства.
\item При встречном движении вместо суммы скоростей используют разность.
\item После вычисления не проверяют, что найденное число отвечает вопросу задачи.
\end{itemize}

\clearpage
\thispagestyle{fancy}
\begin{center}
{\Large\color{deep}Алгоритм: как найти часть или целое}\par
\vspace{5mm}
\end{center}
\begin{center}
\begin{tikzpicture}[
 >={Stealth[length=2.7mm]},
 font=\small,
 start/.style={draw=deep,rounded corners=11pt,minimum width=88mm,minimum height=12mm,align=center,inner sep=5pt},
 decision/.style={draw=accent,diamond,aspect=2.45,align=center,text width=46mm,inner sep=3pt},
 action/.style={draw=accent,rounded corners=4pt,minimum width=66mm,minimum height=17mm,align=center,text width=55mm,inner sep=5pt},
 final/.style={draw=deep,rounded corners=11pt,minimum width=108mm,minimum height=15mm,align=center,text width=98mm,inner sep=4pt},
 arr/.style={->,draw=muted,line width=0.8pt}]
\node[start] (s) at (0,0) {Определите целое $C$, часть $A$ и долю $p$\\($p=\dfrac ab$ или $p=\dfrac{r}{100}$)};
\node[decision] (q) at (0,-3.2) {Что требуется найти?};
\node[action] (left) at (-3.7,-6.6) {\textbf{Часть $A$}\\Умножьте целое на долю\\$A=C\cdot p$};
\node[action] (right) at (3.7,-6.6) {\textbf{Целое $C$}\\Разделите часть на долю\\$C=A:p$, $p\ne0$};
\node[final] (end) at (0,-9.5) {Проверьте результат обратным действием\\и ответьте на вопрос задачи};
\draw[arr] (s) -- (q);
\draw[arr] (q.south west) -- ++(-0.6,-0.4) -| (left.north);
\draw[arr] (q.south east) -- ++(0.6,-0.4) -| (right.north);
\draw[arr] (left.south) -- ++(0,-0.8) -| ([xshift=-32mm]end.north);
\draw[arr] (right.south) -- ++(0,-0.8) -| ([xshift=32mm]end.north);
\end{tikzpicture}
\end{center}

\clearpage
\begin{center}
{\LARGE\color{deep}Тренировочный блок}\par
\vspace{2mm}
{\normalsize Путь А --- 10 заданий}\par
\end{center}
\ruleline
\textit{Решайте последовательно. Записывайте уравнения и указывайте единицы измерения.}
\begin{enumerate}[label=\textbf{\arabic*.},leftmargin=9mm,itemsep=7pt]
\item Найдите $\dfrac38$ числа $64$.
\item Известно, что $\dfrac56$ некоторого числа равны $35$. Найдите это число.
\item Решите уравнение: $4\dfrac12x-4\dfrac12=4\dfrac12$.
\item Вычислите: $0{,}24:0{,}6$.
\item Один из смежных углов в $1\dfrac14$ раза больше другого. Найдите оба угла.
\item Два отряда посадили $1035$ саженцев. Первый посадил $80\%$ от количества, посаженного вторым. Сколько саженцев посадил каждый отряд?
\item Убрали $420$ гектаров --- это $70\%$ всей площади. Найдите площадь поля и площадь, оставшуюся неубранной.
\item После повышения производительности станок начал изготавливать за смену на $84$ детали больше, что составляет $12\%$ прежней выработки. Сколько деталей станок изготавливал раньше?
\item Цена товара снизилась на $18\%$, причём уменьшение составило $72$ рубля. Найдите первоначальную и новую цены.
\item Два автомобиля движутся навстречу друг другу из пунктов, удалённых на $144$ км, и встречаются через $\dfrac45$ часа. Скорость одного автомобиля составляет $80\%$ скорости другого. Найдите скорости обоих автомобилей.
\end{enumerate}

\clearpage
\begin{center}
{\LARGE\color{deep}Ответы}\par
\vspace{3mm}
{\normalsize Таблица к тренировочному блоку}\par
\end{center}
\ruleline
\renewcommand{\arraystretch}{1.4}
\begin{tabularx}{\linewidth}{@{}>{\bfseries}p{14mm}X@{}}
\toprule
№ & Ответ \\
\midrule
1 & $24$ \\
2 & $42$ \\
3 & $x=2$ \\
4 & $0{,}4$ \\
5 & $80^\circ$ и $100^\circ$ \\
6 & Первый отряд --- $460$; второй --- $575$ саженцев \\
7 & Всего $600$ га; осталось $180$ га \\
8 & $700$ деталей \\
9 & Первоначальная цена $400$ руб.; новая цена $328$ руб. \\
10 & $80$ км/ч и $100$ км/ч \\
\bottomrule
\end{tabularx}

\vspace{13mm}
\concept{Итоговая памятка}
\begin{tabular}{@{}>{\raggedright\arraybackslash}p{55mm}>{\raggedright\arraybackslash}p{100mm}@{}}
\toprule
Сигнал в условии & Действие \\
\midrule
Найти долю от целого & Умножить целое на долю \\
Известна доля, найти целое & Разделить известную часть на долю \\
Стало больше на $p\%$ & Новое значение $=C\left(1+\dfrac p{100}\right)$ \\
Стало меньше на $p\%$ & Новое значение $=C\left(1-\dfrac p{100}\right)$ \\
Два объекта едут навстречу & Скорости сложить; затем использовать $S=vt$ \\
\bottomrule
\end{tabular}

\vfill
\begin{center}
{\small\color{muted}Лёвин Артём Александрович эксклюзивно для Надежды Клименко}
\end{center}
\end{document}
